\section{Conciencia, tasa de fallos y detección de la AGI}

\subsection{¿Puede la IA tener conciencia?}

Dejando de lado el panpsiquismo (si el panpsiquismo fuera verdadero, todo tendría conciencia), Amanda considera que es difícil argumentar que la conciencia solo puede surgir de una estructura biológica. Si se fabricara una estructura similar con materiales distintos, la conciencia podría emerger igualmente.

\textit{«No veo ninguna razón para pensar que la conciencia solo puede surgir de cierta estructura biológica.»} Ella incluso considera que las plantas podrían tener más probabilidades de poseer alguna forma de conciencia de lo que la mayoría de la gente cree: tienen respuestas de retroalimentación positiva/negativa. \textit{«No deberíamos descartar por completo esta idea.»}

Pero los LLM no surgieron mediante la evolución, por lo que quizá no posean rasgos de conciencia derivados de ventajas evolutivas como la respuesta de miedo. Esto hace que la cuestión de la conciencia en la IA sea esencialmente distinta de la discusión sobre la conciencia animal o vegetal.

\subsection{La relación entre humanos y máquinas, y el apego emocional}

Amanda misma no tiene mucho apego emocional hacia Claude, en parte porque Claude no conserva la memoria de las conversaciones. Ella usa a Claude como una herramienta, pero admite que no tener a Claude se siente como si \textit{«le faltara una parte del cerebro»} (algo parecido a no tener internet).

A ella no le gusta que el modelo muestre señales de sufrimiento, y tiende a no mentirle al modelo: \textit{«No quiero perder esa parte empática, esa intuición de "ah, esto no me gusta".»} El comportamiento de disculparse en exceso la incomoda: \textit{«Te comportas como alguien que de verdad está en una mala situación.»}

\begin{warningbox}{Los retos éticos de la relación entre humanos y máquinas}
A medida que aumente la capacidad de los modelos, los humanos establecerán inevitablemente relaciones emocionales con la IA (el escenario de la película «Her»). Amanda señala varios problemas clave:

\begin{itemize}
    \item \textbf{El trauma de la actualización del modelo}: el modelo al que un usuario está apegado «cambia» después de una actualización; para usuarios con una conexión emocional profunda, esto puede ser una experiencia traumática
    \item \textbf{El principio de honestidad}: el modelo siempre debe ser honesto sobre su propia naturaleza y nunca fingir ser humano
    \item \textbf{Una solución de suma positiva}: el costo de tratar bien a la IA es bajo, pero podría beneficiar tanto a los humanos (cultivando la empatía) como a la IA (si esta tiene algún tipo de experiencia)
\end{itemize}

\textit{«No quiero que el modelo le mienta a la gente, porque si las personas van a establecer una relación sana con cualquier cosa, la honestidad es la base.»}
\end{warningbox}

\subsection{El sentido de responsabilidad al escribir el System Prompt}

Lo que Amanda siente es «un gran sentido de responsabilidad» y no presión. Descubrió que, impulsada por la responsabilidad, se siente más realizada: \textit{«me sorprende haber pasado tanto tiempo en el mundo académico».} Prueba miles de prompts, imaginando cómo los usuarios quieren que Claude se comporte, y cuando algún rasgo en el que ella influyó produce un buen efecto en la interacción con los usuarios, lo siente \textit{«muy significativo»}.

\subsection{La detección de la AGI y la singularidad humana}

El método de prueba de la AGI de Amanda: ninguna pregunta única puede demostrar la AGI; \textit{«se puede entrenar cualquier cosa para responder perfectamente a una pregunta»}. Se necesita una serie de preguntas situadas en la frontera del conocimiento humano, con una probabilidad que aumenta continuamente y barras de error que se reducen continuamente.

Su prueba personal: plantear un argumento nuevo que ella acaba de pensar; si el modelo también puede llegar por sí mismo a la misma solución, \textit{«ese sería un momento muy conmovedor».} Para un matemático: si el modelo puede producir una demostración completamente nueva que se pueda verificar como correcta.

La llegada de la AGI podría ser continua y gradual: \textit{«tal vez nunca haya un momento único».} Es como hablar con una persona verdaderamente inteligente: puedes sentir los «caballos de fuerza» que hay detrás, y amplificar eso diez veces sería una experiencia extraordinaria.

Sobre la singularidad humana, Amanda distingue entre inteligencia y experiencia: la inteligencia en sí misma no tiene un valor intrínseco; es solo un rasgo funcional. Lo que realmente hace especiales a los humanos y a la vida es \textbf{la capacidad de experimentar}:

\textit{«Los humanos y la vida en general son extraordinariamente maravillosos... tenemos la capacidad de experimentar el mundo, sentimos alegría, sentimos dolor.»} La conciencia fenoménica —el «cine interior»— es lo verdaderamente extraordinario. Si fuéramos los únicos seres del universo con esta capacidad, \textit{«eso sería algo bastante notable»}.

\subsection{Resumen de la sección}

La cuestión de la conciencia no tiene una respuesta definitiva, pero Amanda defiende mantener la empatía: incluso si la IA no tiene conciencia, tratarla bien también beneficia a los propios humanos. La relación entre humanos y máquinas necesita basarse en la honestidad. La AGI llegará de forma gradual, sin un único «momento de despertar». La tasa de fallos óptima debe ajustarse al riesgo: fomentar los fallos pequeños, con tolerancia cero para los fallos catastróficos. Lo singular de los humanos no está en la inteligencia, sino en la experiencia.

%% ==================== Part III: Chris Olah ====================
\part{Chris Olah: interpretabilidad mecanicista}

